\documentclass[10pt,a4paper]{amsart}
\usepackage[margin=19mm]{geometry}
\usepackage{amsmath,amssymb,mathtools}
\usepackage[scr=rsfs,cal=euler]{mathalfa}
\usepackage{xfrac}
\usepackage[unicode,hidelinks,bookmarks=false]{hyperref}
\newcommand{\ie}{i.e.\ }
\newcommand{\cf}{cf.\ }
\newcommand{\resp}{respectively\ }
\newcommand{\Subsection}{\subsection}
\providecommand{\nobreakdash}{\nobreak\hbox{-}}
\setlength{\emergencystretch}{2em}
\multlinegap0pt
\usepackage{bm}
\usepackage{bbm}
\usepackage{dsfont}
\usepackage{xspace}
\usepackage{cancel}
\usepackage{mathtools}
\usepackage[shortlabels]{enumitem}
\usepackage{amsthm}
\makeatletter
\@ifundefined{claim}{\newtheorem{claim}{Claim}}{}
\@ifundefined{sublemma}{\newtheorem{sublemma}[claim]{Lemma}}{}
\makeatother
\title[Pesin theory for transcendental maps and applications]{Pesin theory for transcendental maps and applications}
\author{Anna Jov\'e}
\date{Source: arXiv:2410.19703v2 (CC BY 4.0)}
\begin{document}
\maketitle

\noindent\emph{Source and licence.} Pesin theory for transcendental maps and applications. Version arXiv:2410.19703v2 (CC BY 4.0), distributed under the Creative Commons Attribution 4.0 International licence, as recorded in the arXiv licence metadata for this exact version and shown on the abstract page https://arxiv.org/abs/2410.19703v2. Full rights notice: licenses/arxiv-2410.19703-CC-BY-4.0.txt.\par\smallskip

\noindent\textit{Editorial scope.} Counted unit: the Sublemma whose source label is \texttt{sublemma-puntperiodic1} in the article (source0.tex lines 1086--1095, source label \texttt{sublemma-puntperiodic1}), delivered together with the article's own complete original proof of it (source0.tex lines 1097--1142, one \texttt{proof} environment of 46 lines). No mathematics is written, repaired or re-derived: statement and proof are the article's own text line for line. Required context retained uncounted: none. Every definition, display and label of those lines is retained unchanged. Omitted from this package and not redistributed: 1--1085, 1096--1096, 1143--1447. Nothing inside a retained excerpt is omitted or altered: every retained body is delivered exactly as the article's lines are written, so no omission receipt of any kind accompanies this package. Citations: the retained excerpts carry no citation command at all, so no bibliography entry of the article is transcribed and no reference is redistributed. Reference adaptation: none was needed and none was made; every reference inside the delivered unit resolves inside it, so no unresolved reference or citation is produced. Printed slips kept literal: no printed word, symbol, hypothesis or number of the counted statement or of its proof was altered. Layout and class: the article's own class is \texttt{article} with packages including amsmath, amssymb, graphicx, subcaption, amsfonts, latexsym, amsthm, mathrsfs, mathtools, geometry, url, babel, epsfig, pifont; the wrapper is the lane's pinned \texttt{amsart} editorial wrapper at 10pt on A4, which restates the theorem-like environments the retained text needs and adds the article's own macro definitions that the retained text needs (none), with extra wrapper packages (bm, bbm, dsfont, xspace, cancel, mathtools, enumitem). The delivered numbering is the unit's own. Assets: no retained span contains a figure environment, a graphics-inclusion command, a PDF-page inclusion, a table or a TikZ picture; no image file is copied into this package, the package declares no asset dependency and no image file is redistributed with it. Licence: version arXiv:2410.19703v2 of this article is distributed under the Creative Commons Attribution 4.0 International licence, as recorded in the arXiv licence metadata for this exact version and shown on the abstract page https://arxiv.org/abs/2410.19703v2. A full rights notice is delivered at licenses/arxiv-2410.19703-CC-BY-4.0.txt. Characters: every retained excerpt is pure ASCII, so the delivered engine, XeLaTeX with its default Latin Modern text font, reports no missing character.
\begin{sublemma}\label{sublemma-puntperiodic1}
	There exists a backward orbit $ \left\lbrace \xi_n\right\rbrace _n \subset \partial\mathbb{D}$, and constants $ m\in\mathbb{N} $, $ 0<\rho_m\leq  \rho $, and $ r\in (0, R/2) $  such that:
	\begin{enumerate}[label={\em (1.\arabic*)}]
		\item\label{item-1} $ x_0\coloneqq \varphi^* (\xi_0)$ and $ x_m\coloneqq \varphi^* (\xi_m)$ are well-defined, and $ x_0\in D(x, R/2) $ and $ x_m\in D(x_0, r/3) $;
		\item\label{item-2}  the inverse branch $ F_m $ of $ f^m $ sending $ x_0 $ to $ x_m $ is well-defined in $ D(x_0, r) $, and 
$ \textrm{\em diam } F_m(D(x_0, r))<r/3 $;
\item \label{item-3} the inverse branch $ G_m $ of $ g^m $ sending $ \xi_0 $ to $ \xi_m $ is well-defined in $ D(\xi_0, \rho_m) $, and satisfies \[  G_m(R_{\rho_m}(\xi_0))\subset \Delta_{ \rho_m}(\xi_m);\] 
	\item \label{item-4} $ \Delta_{\rho}(\xi_0)\cap  \Delta_{ \rho}(\xi_m)\neq\emptyset$, and, if $ z\in\Delta_{ \rho} (\xi_0)\cup \Delta_{ \rho} (\xi_m)$,  then $ \varphi(z)\in D(x_0, r) $.
	\end{enumerate}
\end{sublemma}	

\begin{proof}
Let $ A_n= D(x^n, r_n) $ be a countable basis for $ D(x, R) $ with the Euclidean topology, where $ x^n\in\partial U $ and $ A_n\subset D(x,R) $. 

In order to apply the hyptothesis of the theorem, we shall construct an appropriate countable sequence of measurable sets $ \left\lbrace K_k\right\rbrace _k $ of $ \partial\mathbb{D} $. We do it as follows.

For all $ n\geq 0 $, let 
\[ K^{n}=\left\lbrace \xi\in\partial\mathbb{D}\colon \varphi^*(\xi)\in D(x^n, {r_n}/{2}) \right\rbrace . \] It is clear that $ \lambda(K^{n})>0 $. By the Lehto-Virtanen Theorem, the angular limit exists whenever the radial limit exists. Therefore, there exists $ \rho_n>0 $ small enough so that
\[ K^n_{\rho_n}=\left\lbrace \xi\in K^n\colon \Delta_{\rho_n}(\xi)\subset D(x^n, {r_n}/{2}) \right\rbrace  \]
has positive $ \lambda $-measure. We can assume that every point in $ K^n_{\rho_n} $  is a Lebesgue density point  for $ K^n_{\rho_n} $. Then, if we take
$ \xi^n \in  K^n_{\rho_n} $, there exists a circular interval  $ I_{\xi^n} $ around $ \xi^n $ such that for any $ \zeta_1, \zeta_2\in I_{\xi^n}  $,
\[\Delta_{\rho_n}(\zeta_1)\cap  \Delta_{ \rho_n}(\zeta_2)\neq\emptyset.\] Then, $ K^n_{\rho_n} \cap  I_{\xi^n} $ has positive $ \lambda $-measure.
Note that this property only depends on the length of the interval, as long as $  \xi^n $ is a Lebesgue density point for $ K^n_{\rho_n} $.
Then, it is clear that there exist finitely many circular intervals $ I_1^n, \dots, I^n_{i_n} $ with this property. 

Let \[K^{1,n}_{*, i} \coloneqq K^n_{\rho_n} \cap  I^n_{i},\hspace{0.5cm} i=1, \dots, i_n,\]
\[
K^{1,n}_*\coloneqq\left\lbrace K^{1,n}_{*,1},\dots, K^{1,n}_{*,i_n}\right\rbrace.
\]
Then, we define the set $ K^{j,n}_*$, as before, but replacing $ \rho_n $ by $ \rho_n /2^j$.

Having introduced all this notation of the sets $ \left\lbrace K^{j,n}_* \right\rbrace _{n,j}$, we arrange the sequence $ \left\lbrace K_k\right\rbrace _k $ as follows. We construct this sequence of sets inductively, adding at each step finitely many sets.
Indeed, let us start by putting the block $ K^{1,1}_*\coloneqq\left\lbrace K^{1,1}_{*,1},\dots, K^{1,1}_{*,i_1}\right\rbrace $ as the first elements of the sequence. Then,
for the $ k $-th step of the induction, we consider $ A_{k} $ and let $ A_{k_1}, \dots,  A_{k_n}$ be all the sets of $ A_1,\dots, A_n $ such that $ A_n\subset A_{k_i} $. Then, we add to the sequence the blocks
\[K^{1,k_1}_*, \dots K^{1,k_n}_*, \dots, K^{k,k_1}_*, \dots, K^{k,k_n}_*.\]


 Basically, the idea is that, when one set is in the sequence $ \left\lbrace K_k\right\rbrace _k $ for the first time, then it appears infinitely often. Moreover, the set of points in $ \left\lbrace K_k\right\rbrace _k $ has  measure $ \lambda ((\varphi^*)^{-1}(D(x,R))) $. Indeed, the set of points in $ \partial\mathbb{D} $ for which the radial limit exists has full measure. Let $ \zeta $ be one of such points. Then, $ \varphi^*(\zeta) \in A_j$, for some $ j $, and for $ \rho>0 $ small enough, $ \Delta_\rho (\zeta)\subset A_j $. Then, there exists $ n\geq 0 $ such that $ A_n\subset A_j $ and $ \rho<\rho_j/2^n $, so $ \zeta \in K^n_{*, k_n}$, as desired.

By the assumption of the theorem, for $ \lambda $-almost every $ \xi_0\in\partial\mathbb{D} $, 
there exists a backward orbit $ \left\lbrace \xi_n\right\rbrace _n $ such that  the hypothesis on the definition of the inverse branches for $ \left\lbrace \xi_n\right\rbrace _n $ and $ \left\lbrace x_n\coloneqq \varphi^*(\xi_n)\right\rbrace _n $ are accomplished, and there exists $ n_k\to \infty  $ with $ \xi_{n_k}\in K_k $. 

Without loss of generality, we assume $ \xi_0 $ is chosen so that $ x_0\in D(x, R/2) $. Let $ r>0 $ be such that the inverse branches realizing the backward orbit $ \left\lbrace x_n\right\rbrace _n $ are well-defined in $ D(x_0,r) $. There is no loss of generality on assuming $ r\in (0, R/2) $. 

On the one hand, since $ \left\lbrace A_n \right\rbrace _n$ is a basis for $ D(x,R) $, there exists $ n_0 $ such that $$ x_0\in A_{n_0}\subset D(x_0, r/3) , $$ and  $ \xi_0 \in K^{n_0}_*$,  by the previous remark. In particular, for $ \rho_{n_0} $,  
\[\Delta_{\rho_{n_0}}(\xi_0)\subset A_{n_0}\subset D(x_0, r/3).\]
On the other hand, 
by the construction of the sets $ \left\lbrace K_n\right\rbrace _n  $, the backward orbit visits $ D(x_0,r) $ infinitely many times. Let $ n_1 $ be large enough so that, for all $ n\geq n_1 $, if $ x_n\in D(x_0, r) $, then $ \textrm{diam }F_n(D(x_0, r))<r/3  $. 

 By the construction of the sets $ \left\lbrace K_n\right\rbrace _n  $, there exists $ m\geq \max \left\lbrace n_0,  n_1\right\rbrace  $ such that $ \xi_m \in K^{n_0}_*$. Hence, we  take $ r>0 $, $ \rho=\rho_{n_0} $, and $ \xi_0 $ and $ \xi_m $ as above, and define $ \rho_m>0 $ as the radius such that the inverse branch $ G_m $ sending $ \xi_0 $ to $ x_m $ is defined around $ \xi_0 $ (such a radius exists by our assumptions on the orbit $ \left\lbrace \xi_n\right\rbrace _n $). We have to check that, with these choices, the requirements are satisfied. 
 

First, by the choice of $ \left\lbrace \zeta_n\right\rbrace _n $, $ \varphi^*(\xi_0)\eqqcolon x_0 $ and $ \varphi^*(\xi_m)\eqqcolon x_m $ are well-defined. Moreover, by the choice of $ r $, we have $ x_0\in D(x, R/2) $. Since $ \xi_0, \xi_m \in K^{n_0}_*$, we have 
\[\Delta_{\rho}(\xi_0)\cap  \Delta_{ \rho}(\xi_n)\neq\emptyset,\]
and $ \Delta_{\rho}(\xi_0), \Delta_{\rho}(\xi_m)\subset A_{n_0}\subset D(x_0, r/3) $. In particular, $ x_m\in D(x_0, r/3) $, so (1.1) and (1.4) hold. To see (1.2), note that $ r $ has been chosen so that the inverse branches corresponding to $ \left\lbrace x_n\right\rbrace _n $ are well-defined in $ D(x_0,r) $, and $ m $ is large enough so that $ \textrm{diam }F_m(D(x_0, r))<r/3  $, as desired.
 Requirement (1.3)  is directly satisfied by the choice of $ \rho_m $. Therefore, we have proved the lemma.
\end{proof}

\end{document}
